\documentclass[10pt,a4paper]{amsart}
\usepackage[margin=19mm]{geometry}
\usepackage{amsmath,amssymb,mathtools}
\usepackage[scr=rsfs,cal=euler]{mathalfa}
\usepackage{xfrac}
\usepackage[unicode,hidelinks,bookmarks=false]{hyperref}
\newcommand{\ie}{i.e.\ }
\newcommand{\cf}{cf.\ }
\newcommand{\resp}{respectively\ }
\newcommand{\Subsection}{\subsection}
\providecommand{\nobreakdash}{\nobreak\hbox{-}}
\setlength{\emergencystretch}{2em}
\multlinegap0pt
\usepackage{bm}
\usepackage{bbm}
\usepackage{dsfont}
\usepackage{xspace}
\usepackage{cancel}
\usepackage{mathtools}
\usepackage{amsthm}
\makeatletter
\@ifundefined{theorem}{\newtheorem{theorem}{Theorem}}{}
\@ifundefined{assumption}{\newtheorem{assumption}{Assumption}}{}
\@ifundefined{lemma}{\newtheorem{lemma}[theorem]{Lemma}}{}
\@ifundefined{proposition}{\newtheorem{proposition}{Proposition}}{}
\makeatother
\newcommand{\N}{\mathcal N}
\newcommand{\lr}[1]{\left(#1\right)}
\newcommand{\norm}[2]{\left\|#1\right\|_{#2}}
\newcommand{\pr}{\pi_{\overline{\range(A)}}}
\newcommand{\range}{\mathcal R}
\newcommand{\scalar}[2]{\langle{ #1},{#2} \rangle}
\title[Regularization of the metric generalized inverse in Banach s]{Regularization of the metric generalized inverse in Banach spaces and the dichotomy phenomenon}
\author{Peter Math\'e, Bernd Hofmann}
\date{Source: arXiv:2606.26730v1 (CC BY 4.0)}
\begin{document}
\maketitle

\noindent\emph{Source and licence.} Regularization of the metric generalized inverse in Banach spaces and the dichotomy phenomenon. Version arXiv:2606.26730v1 (CC BY 4.0), distributed under the Creative Commons Attribution 4.0 International licence, as recorded in the arXiv licence metadata for this exact version and shown on the abstract page https://arxiv.org/abs/2606.26730v1. Full rights notice: licenses/arxiv-2606.26730-CC-BY-4.0.txt.\par\smallskip

\noindent\textit{Editorial scope.} Counted unit: the Lemma whose source label is \texttt{lem:B} in the article (source0.tex lines 478--485, source label \texttt{lem:B}), delivered together with the article's own complete original proof of it (source0.tex lines 486--525, one \texttt{proof} environment of 40 lines). No mathematics is written, repaired or re-derived: statement and proof are the article's own text line for line. Required context retained uncounted: ass:geometry (lines 115--117); eq:9 (lines 464--467); prop:decomposition (lines 874--885). Every definition, display and label of those lines is retained unchanged. Omitted from this package and not redistributed: 1--114, 118--463, 468--477, 526--873, 886--1167. Nothing inside a retained excerpt is omitted or altered: every retained body is delivered exactly as the article's lines are written, so no omission receipt of any kind accompanies this package. Citations: the retained excerpts carry no citation command at all, so no bibliography entry of the article is transcribed and no reference is redistributed. Reference adaptation: none was needed and none was made; every reference inside the delivered unit resolves inside it, so no unresolved reference or citation is produced. Printed slips kept literal: no printed word, symbol, hypothesis or number of the counted statement or of its proof was altered. Layout and class: the article's own class is \texttt{amsart} with packages including inputenc, xcolor; the wrapper is the lane's pinned \texttt{amsart} editorial wrapper at 10pt on A4, which restates the theorem-like environments the retained text needs and adds the article's own macro definitions that the retained text needs (\texttt{\textbackslash N}, \texttt{\textbackslash lr}, \texttt{\textbackslash norm}, \texttt{\textbackslash pr}, \texttt{\textbackslash range}, \texttt{\textbackslash scalar}), with extra wrapper packages (bm, bbm, dsfont, xspace, cancel, mathtools). The delivered numbering is the unit's own. Assets: no retained span contains a figure environment, a graphics-inclusion command, a PDF-page inclusion, a table or a TikZ picture; no image file is copied into this package, the package declares no asset dependency and no image file is redistributed with it. Licence: version arXiv:2606.26730v1 of this article is distributed under the Creative Commons Attribution 4.0 International licence, as recorded in the arXiv licence metadata for this exact version and shown on the abstract page https://arxiv.org/abs/2606.26730v1. A full rights notice is delivered at licenses/arxiv-2606.26730-CC-BY-4.0.txt. Characters: every retained excerpt is pure ASCII, so the delivered engine, XeLaTeX with its default Latin Modern text font, reports no missing character.
\begin{assumption}[Banach space geometry]\label{ass:geometry}
We consider a bounded linear operator $A \colon X \to Y$, where both spaces~$X$ and~$Y$ are \emph{reflexive} and \emph{strictly convex} Banach spaces.
\end{assumption}

\begin{equation}
  \label{eq:9}
  B:= \lr{J_{X}}^{-1} A^{\prime} J_{Y}\colon Y \to X,
\end{equation}

\begin{proposition}\label{prop:decomposition}
  Under Assumption~\ref{ass:geometry} the following holds true.
  % \label{it:decomposition}
  Given the linear operator~$A\colon X \to
  Y$, we have the decompositions
  \begin{align}
    X &= \N(A) + J_{X}^{-1}\lr{\N(A)^{\perp}}\label{it:Xdec}
    \intertext{and}
    Y&= \overline{\range(A)} + J_{Y}^{-1}\lr{\range(A)^{\perp}}.\label{it:Ydec}
  \end{align}
  The mapping~$A\colon  J_{X}^{-1}\lr{\N(A)^{\perp}} \to \range(A)$ is injective.
\end{proposition}

\begin{lemma}\label{lem:B}
  Let~$B$ be as in~(\ref{eq:9}). The following holds true.
  \begin{enumerate}
  \item\label{it:Bpi} We have that~$\pi_{\N(A)}B y=0,\ y\in Y$,
    \item \label{it:piB}$B \pr y = B y,\ y\in Y$.
  \item An element~$y\in Y$ fulfills~$A B y=0$ if and only if~$\pr y=0$.
  \end{enumerate}
\end{lemma}

\begin{proof}
  In the proof we use Proposition~\ref{prop:decomposition}.
  To prove item~\eqref{it:Bpi} it is required to show that~$B\lr{y -
    \pi_{\overline{\range(A)}}y} = 0$. Let $y\in C_Y(A)=
  J_{Y}^{-1}\lr{\range(A)^{\perp}}$,  we have to show that~$B y=0$.
 We notice that~$\range(A)^{\perp} = \N(A^{\prime})$, and hence there
  is~$y^{\prime}\in \N(A^{\prime})$ with~$y= J_{Y}^{-1}(y^{\prime})$.
  But
  $$
 B y = J_{X}^{-1} A^{\prime} J_{Y}(y) = J_{X}^{-1}
 A^{\prime} J_{Y}J_{Y}^{-1}(y^{\prime}) = J_{X}^{-1} A^{\prime}y^{\prime}= 0,
 $$
 which proves the first assertion.
 \par
 Similarly we proceed for item~\eqref{it:piB}. Due to the
 decomposition of~$Y$ it requires to show that~$\range(B)\subset
 J_{X}^{-1}\lr{\N(A)^{\perp}} =
 J_{X}^{-1}\lr{\overline{\range(A^{\prime})}}$.
 This is immediate from the definition of~$B$ in~(\ref{eq:9}).
 \par
 It remains to prove the last equivalence.
 Let us first suppose that \newline $\pr y =0$. By using the decomposition of~$Y$ this
 means that~$y \in J_{Y}^{-1}\lr{\overline{\range(A)}^{\perp}}$, which is
 equivalent to~$y\in J_{Y}^{-1}\lr{\N(A^{\prime})}$, and
 henceforth~$J_{Y}y \in \N(A^{\prime})$. It follows
 that~$A^{\prime}J_{Y}y=0$, which yields~$By=0$.
 \par
 To prove the converse, let~$y$ be such that~$ABy=0$. This means
 that~$z:= J_{X}^{-1}A^{\prime}J_{Y}y \in\N(A)$. By definition we have
 that
 $$
\scalar{z }{A^{\prime}J_{Y}y} = \norm{z}{X}^{2}
\norm{A^{\prime}J_{Y}y}{X^{\prime}}^{2} = \scalar{ A z }{J_{Y}y} =0.
$$
Therefore, it must hold true that~$A^{\prime}J_{Y}y=0$, i.e.,\
$J_{Y}y\in \N(A^{\prime}) = \overline{\range(A)}^{\perp}$. We see
that~$y\in J_{Y}^{-1}\lr{\N(A^{\prime})}=J_{Y}^{-1}\lr{
  \overline{\range(A)}^{\perp}} $, and hence in the decomposition
of~$Y$ its projection~$\pr y$ equals zero. The proof is complete.
\end{proof}

\end{document}
